\documentclass[11pt,a4paper]{article}

\usepackage[utf8]{inputenc}
\usepackage[T1]{fontenc}
\usepackage[french]{babel}
\usepackage{lmodern}
\usepackage[a4paper,margin=2.2cm]{geometry}
\usepackage{amsmath,amssymb,amsthm}
\usepackage{mathrsfs}
\usepackage{fancyhdr}
\usepackage{lastpage}
\usepackage{xcolor}
\usepackage{enumitem}
\usepackage{tabularx}
\usepackage{array}
\usepackage{titlesec}
\usepackage{framed}
\usepackage{tcolorbox}
\usepackage{hyperref}
\usepackage{setspace}

\hypersetup{
    colorlinks=true,
    linkcolor=black,
    urlcolor=blue
}

\setlength{\parindent}{0pt}
\setlength{\parskip}{0.35em}
\setstretch{1.08}

\definecolor{bleufonce}{RGB}{25,45,85}
\definecolor{bleuclair}{RGB}{235,242,252}
\definecolor{grisclair}{RGB}{245,245,245}

\pagestyle{fancy}
\fancyhf{}
\fancyhead[L]{\textsc{Terminale spé -- Spécialité mathématiques + Maths expertes}}
\fancyhead[R]{\textsc{Devoir maison 0}}
\fancyfoot[C]{Page \thepage\ /\pageref{LastPage}}

\titleformat{\section}
{\normalfont\Large\bfseries\color{bleufonce}}
{}{0em}{}

\titleformat{\subsection}
{\normalfont\large\bfseries\color{bleufonce}}
{}{0em}{}

\newtcolorbox{enteteDM}{
    colback=bleuclair,
    colframe=bleufonce,
    boxrule=0.8pt,
    arc=2mm,
    left=2mm,
    right=2mm,
    top=2mm,
    bottom=2mm
}

\newtcolorbox{consignesbox}{
    colback=grisclair,
    colframe=black!60,
    boxrule=0.5pt,
    arc=1.5mm,
    left=2mm,
    right=2mm,
    top=1.5mm,
    bottom=1.5mm
}

\newcommand{\R}{\mathbb{R}}
\newcommand{\N}{\mathbb{N}}

\begin{document}

\begin{center}
\begin{enteteDM}
{\LARGE \textbf{DM 0 de mathématiques}}\\[0.3em]
{\large Niveau Terminale générale}\\
{\normalsize Spécialité mathématiques + option mathématiques expertes}\\[0.4em]
\textbf{Durée conseillée :} 2 h 30 à 3 h 30
\end{enteteDM}
\end{center}

\vspace{0.4em}

\begin{consignesbox}
\textbf{Consignes :}
\begin{itemize}[leftmargin=1.2cm]
    \item La calculatrice est interdite.
    \item La rédaction devra être \textbf{claire, rigoureuse et soignée}.
    \item Toute affirmation devra être \textbf{justifiée}.
    \item Note le tend que t'as mis à le faire à la fin de ta copie (chronomètre toi).
    \item La qualité de la présentation, de l'enchaînement logique des arguments et de la rédaction sera prise en compte.
    \item comme d'hab pas de chat gpt, pas d'aide extérieure si tu bloques plus de 25 min sautes question et reviens-y plus tard.
\end{itemize}
\end{consignesbox}

\vspace{0.6em}

\section*{Exercice 1 -- Étude d'une fonction \hfill \textnormal{(9 points)}}

On considère la fonction \( f \) définie sur \( ]0;+\infty[ \) par
\[
f(x)=x+\frac{1}{x}-2\ln(x).
\]

\subsection*{Partie A -- Étude préliminaire}
\begin{enumerate}[label=\textbf{\arabic*.}, leftmargin=1.1cm]
    \item Déterminer la limite de \( f(x) \) lorsque \( x\to 0^+ \).
    \item Déterminer la limite de \( f(x) \) lorsque \( x\to +\infty \).
    \item Montrer que, pour tout réel \( x>0 \),
    \[
    f'(x)=\frac{(x-1)^2}{x^2}.
    \]
    \item Dresser le tableau de variations de \( f \) sur \( ]0;+\infty[ \).
    \item En déduire que, pour tout réel \( x>0 \),
    \[
    x+\frac{1}{x}\ge 2\ln(x)+2.
    \]
\end{enumerate}

\subsection*{Partie B -- Exploitation}
On considère l'équation
\[
x+\frac{1}{x}-2\ln(x)=3.
\]

\begin{enumerate}[label=\textbf{\arabic*.}, start=6, leftmargin=1.1cm]
    \item Justifier que cette équation admet exactement deux solutions sur \( ]0;+\infty[ \).
    \item Montrer que l'une de ces solutions appartient à \( ]0;1[ \) et l'autre à \( ]1;+\infty[ \).
    \item Donner un encadrement d'amplitude \(10^{-2}\) de chacune de ces deux solutions.
\end{enumerate}

\vspace{0.5em}

\section*{Exercice 2 -- Fonction et suite récurrente \hfill \textnormal{(14 points)}}

On considère la fonction \( g \) définie sur \( [0;+\infty[ \) par
\[
g(x)=\frac{2x+3}{x+2}.
\]
On considère également la suite \( (u_n) \) définie par
\[
u_0=2
\qquad \text{et} \qquad
u_{n+1}=g(u_n)
\quad \text{pour tout entier naturel } n.
\]

\subsection*{Partie A -- Étude de la fonction}
\begin{enumerate}[label=\textbf{\arabic*.}, leftmargin=1.1cm]
    \item Déterminer \( g(0) \).
    \item Montrer que, pour tout réel \( x\ge 0 \),
    \[
    g(x)=2-\frac{1}{x+2}.
    \]
    \item En déduire la limite de \( g(x) \) lorsque \( x\to +\infty \).
    \item Montrer que, pour tout réel \( x\ge 0 \),
    \[
    g'(x)=\frac{1}{(x+2)^2}.
    \]
    \item Étudier le sens de variation de \( g \) sur \( [0;+\infty[ \).
    \item Résoudre sur \( [0;+\infty[ \) l'équation
    \[
    g(x)=x.
    \]
\end{enumerate}

\subsection*{Partie B -- Étude de la suite}
\begin{enumerate}[label=\textbf{\arabic*.}, start=7, leftmargin=1.1cm]
    \item Calculer \( u_1 \) puis \( u_2 \).
    \item Montrer par récurrence que, pour tout entier naturel \( n \),
    \[
    u_n\ge 1.
    \]
    \item Montrer que, pour tout entier naturel \( n \),
    \[
    u_{n+1}-1=\frac{u_n+1}{u_n+2}.
    \]
    \item Montrer que, pour tout entier naturel \( n \),
    \[
    u_{n+1}-u_n=\frac{(1-u_n)(u_n+1)}{u_n+2}.
    \]
    \item En déduire le sens de variation de la suite \( (u_n) \).
    \item Montrer que la suite \( (u_n) \) converge.
    \item Déterminer sa limite.
\end{enumerate}

\subsection*{Partie C -- Interprétation graphique}
Dans un repère, on note \( \mathcal C_g \) la courbe représentative de la fonction \( g \), et \( \Delta \) la droite d'équation \( y=x \).

\begin{enumerate}[label=\textbf{\arabic*.}, start=14, leftmargin=1.1cm]
    \item Expliquer comment la position relative de \( \mathcal C_g \) et de \( \Delta \) permet de retrouver le sens de variation de la suite \( (u_n) \).
\end{enumerate}

\vspace{0.5em}

\section*{Exercice 3 -- Taux de variation et tangente \hfill \textnormal{(7 points)}}

On considère la fonction \( h \) définie sur \( \R \) par
\[
h(x)=x\sqrt{x^2+1}.
\]

\begin{enumerate}[label=\textbf{\arabic*.}, leftmargin=1.1cm]
    \item Calculer \( h(0) \).
    \item Étudier la parité de la fonction \( h \).
    \item Montrer que, pour tout réel \( x \),
    \[
    h'(x)=\frac{2x^2+1}{\sqrt{x^2+1}}.
    \]
    \item Déterminer une équation de la tangente à la courbe représentative de \( h \) au point d'abscisse \( 1 \).
    \item Déterminer, s'il existe, un point de la courbe en lequel la tangente est parallèle à la droite d'équation
    \[
    y=3x-2.
    \]
    \item Étudier la dérivabilité de \( h \) en \( 0 \).
    \item En déduire l'équation de la tangente à la courbe au point d'abscisse \( 0 \).
\end{enumerate}

\vspace{0.5em}

\section*{Exercice 4 -- Probabilités conditionnelles \hfill \textnormal{(7 points)}}

Une urne \( A \) contient 3 boules rouges et 2 boules bleues.\\
Une urne \( B \) contient 1 boule rouge et 4 boules bleues.

On choisit au hasard une urne, avec équiprobabilité, puis on tire une boule dans l'urne choisie.

On note :
\begin{itemize}[leftmargin=1.2cm]
    \item \( U_A \) l'événement : « l'urne choisie est l'urne \( A \) » ;
    \item \( R \) l'événement : « la boule tirée est rouge ».
\end{itemize}

\begin{enumerate}[label=\textbf{\arabic*.}, leftmargin=1.1cm]
    \item Calculer \( P(U_A) \) et \( P(U_B) \).
    \item Calculer \( P_{U_A}(R) \) et \( P_{U_B}(R) \).
    \item En déduire \( P(R) \).
    \item Calculer \( P_R(U_A) \).
    \item On répète l'expérience 3 fois, de manière indépendante. On note \( X \) le nombre de fois où l'on obtient une boule rouge.
    \begin{enumerate}[label=\textbf{\alph*)}, leftmargin=0.9cm]
        \item Justifier que \( X \) suit une loi binomiale dont on précisera les paramètres.
        \item Calculer \( P(X=2) \).
        \item Calculer \( P(X\ge 1) \).
    \end{enumerate}
\end{enumerate}

\vspace{0.5em}

\section*{Bonus -- Question ouverte \hfill \textnormal{(3 points)}}

Déterminer tous les réels \( x>0 \) tels que
\[
x+\frac{1}{x}=2+2\ln(x).
\]

\vspace{1em}

\begin{center}
\begin{tcolorbox}[colback=bleuclair,colframe=bleufonce,boxrule=0.7pt,arc=2mm,width=0.95\textwidth]
\centering
\textbf{\Large Barème indicatif}
\end{tcolorbox}
\end{center}

\vspace{-0.3em}

\renewcommand{\arraystretch}{1.35}
\begin{center}
\begin{tabular}{|>{\centering\arraybackslash}m{8cm}|>{\centering\arraybackslash}m{3cm}|}
\hline

\textbf{Exercice} & \textbf{Points} \\
\hline
Exercice 1 -- Étude d'une fonction & 9 \\
\hline
Exercice 2 -- Fonction et suite récurrente & 14 \\
\hline
Exercice 3 -- Taux de variation et tangente & 7 \\
\hline
Exercice 4 -- Probabilités conditionnelles & 7 \\
\hline
Bonus -- Question ouverte & 3 \\
\hline

\textbf{Total principal} & \textbf{37} \\
\hline
\end{tabular}
\end{center}

\vspace{0.8em}

\begin{consignesbox}
\textbf{Appréciation qualitative attendue :}
\begin{itemize}[leftmargin=1.2cm]
    \item qualité de la rédaction ;
    \item rigueur des démonstrations ;
    \item enchaînement logique des arguments ;
    \item soin apporté à la copie ;
    \item capacité à mobiliser les outils du cours de manière pertinente.
\end{itemize}
\end{consignesbox}

\end{document}